\section{轨迹预处理与评价方法}
\label{sec:method}

\subsection{何时应将相邻观测断开}
采样空档两端的观测只能说明两个时刻的位置，不能说明中间走过的路径。因此，分段首先处理不可靠的连接，而不是立即判断某个端点是否错误。

记第$i$个观测的工作平面坐标为$p_i=(E_i,N_i)$，时间为$t_i$。当前相邻观测满足
\[
\Delta t_i=t_{i+1}-t_i>\tau_t,\qquad
\Delta t_i<0,\qquad\text{或}\qquad
\|p_{i+1}-p_i\|>\tau_d
\]
中的任一条件时，在右侧点之前断开。等于阈值时不切分，右点只归入新片段。原始数据未出现负时间差，这一情况仍保留为输入检查。时间差为零不单独触发时间切分，也不被填成可计算的速度。

切开后的短片段是否还值得继续处理，是另一个问题。参考设置过滤点数少于$n_{\min}$或累计长度小于$L_{\min}$的片段；累计长度只加总段内相邻边，不包含已经切断的连接。将切分和过滤分开，才能辨认信息是在断点处失去连接，还是在过滤时被整段舍弃。下文用S表示这两个连续操作。

参考条件是否合适，需要后面的参数实验判断。这里首先核对每个原始点的片段归属、边界位置和过滤原因，并检查输出是否重新跨过原始断点。相同点不能因切分重复出现在两个片段，过滤掉的记录也继续计入结果分母。

\subsection{怎样检查方向突变，而不将它直接当作噪声}
方向突然改变可能来自位置偏移，也可能是真实折返。实验以课程给定的双侧方向关系\cite{teacher}识别待检查点，而不把方向变化预先视为噪声标签。

对非零位移边，定义当前点到下一点的出边方向：
\[
\alpha_i=\left[\frac{180}{\pi}\operatorname{atan2}(E_{i+1}-E_i,N_{i+1}-N_i)\right]\bmod 360^\circ.
\]
方向差采用环形角差$\delta(a,b)=\min(|a-b|,360^\circ-|a-b|)$。若
\[
\delta(\alpha_i,\alpha_{i-1})>\theta
\quad\text{且}\quad
\delta(\alpha_i,\alpha_{i+1})>\theta,
\]
则标记点$i$。这里比较相邻出边方向，需要相应的前后观测；它不是把任意三点的一个内角与阈值比较。

若边遍历边删除，前面的动作会改变后续邻域，结果可能依赖遍历次序。实验先在本阶段完整输入上计算全部标记，再同时删除，不迭代至不再能删。首尾、缺少必要后继或零位移导致方向不可计算的窗口保留并标明原因。此后重新计算相邻关系，避免把删除前的方向和长度继续用于下一步。

核验时逐个检查实际删除点是否满足这次输入上的谓词。构造折返和尖峰例子则用于检查规则的解释边界：\textbf{正确执行方向规则，与准确辨认真实噪声不是同一个判断。}这一步记为D。

\subsection{怎样减少点数并控制简化偏差}
短时间内的重复或近共线观测会增加存储量，但任意删点又可能拉直转弯。DP的做法是先用端点连线近似一段轨迹，若中间点偏离过大，就保留最远点，并分别处理左右两段\cite{teacher}。

对于端点$a,b$与中间点$p$，投影须落在有限线段内：
\[
u=\operatorname{clip}_{[0,1]}\!\left(\frac{(p-a)\cdot(b-a)}{\|b-a\|^2}\right),
\qquad d(p,[a,b])=\|p-a-u(b-a)\|.
\]
端点重合时取点到端点的距离。最大距离严格大于$\varepsilon$才继续拆分；等于容差时保留端点。实验另外将$\varepsilon=0$定义为不简化、保留全部输入的参照模式。

输出保留原始索引，而不靠浮点坐标反查对应关系。每条输出线段都对应一段完整输入，独立评价逐区间重算偏差，检查端点、顺序和容差。普通、重复坐标、回头与退化线段均有明确处理。

这一步记为P。\textbf{P的误差保证针对进入P的完整输入。}在D中已经删除的点不属于该输入；若P之后再执行其他删除，其即时误差也不能自动代表最后输出。因此还需要下面的完整处理链评价。

\subsection{为什么不能只比较删点率或最后一步误差}
一种方案可能先过滤掉困难片段，再在剩余点上得到很小误差。若只看候选自己的结果，就可能把信息丢失误认为处理改进。为避免这个问题，实验在任何候选处理前，用原始30秒与400工作米条件建立共同窗口，不预先过滤短段，也不先去噪。

对方法$M$，原始点自己被保留，或者它被同一最终片段中的相邻输出索引夹住、且属于同一共同窗口时，计为覆盖。相应偏差是该点到夹持线段的有限线段距离。没有这样的表示关系时记为未覆盖，不从其他窗口寻找最近线段，也不向片段两端外推。记覆盖集合为$C(M)$，每个已覆盖点的偏差为$e_M(i)$。

\begin{table}[htbp]\centering\small
\caption{互补的评价读数及其参考对象}\label{tab:metrics}
\begin{tabularx}{\linewidth}{@{}L{28mm}L{55mm}Y@{}}\toprule
读数 & 计算对象 & 回答的问题\\\midrule
显式保留率 & $N_{\mathrm{final}}/N_{\mathrm{raw}}$ & 最终实际存储多少原始点\\
共同覆盖率 & $|C(M)|/N_{\mathrm{raw}}$ & 原始范围还有多少得到表示\\
共同几何偏差 & 原始覆盖点到对应最终线段的距离 & 完整处理怎样改变原始折线\\
DP省点率 & $1-N_{P,\mathrm{out}}/N_{P,\mathrm{in}}$ & 简化这一步减少多少点\\
DP即时误差 & 完整P输入到对应P输出区间 & 简化是否满足其容差\\
断点与空输出 & 跨共同断点边、无输出记录等 & 是否丢失连接边界或整条表达\\\bottomrule
\end{tabularx}
\tabnote{覆盖点不一定被实际存储。空输入、无可计算参考时保留不可用状态，不把误差或省点率填成零；几何偏差不是到真实位置或道路的误差。}
\end{table}

S、D及其组合与参考方案$R$比较时，先要求原有覆盖集合不丢失，即$C(R)\subseteq C(M)$，并在同一参考覆盖集合上核对
\[
\max_{i\in C(R)}e_M(i)\leq\max_{i\in C(R)}e_R(i)+\eta,
\]
其中$\eta$只处理浮点数计算的微小差别，不作为允许真实退化的额度。此外，不丢原有非空记录和窗口、不跨原始断点，并满足实际P的共同误差预算。\textbf{这里保护的是每条记录的最大偏差，不是每个点的误差都不增加。}新增覆盖点的偏差另行报告。

全部保护条件通过后，覆盖增加至少一个点，或可比较的最大偏差严格降低，才认为该记录有相应改善。整组自动替换要求所有记录满足保护，且至少一条严格改善；不以其他记录的收益抵消失败。当前方案与最初参考分别比较，避免后续迭代丢回已有收益。

P专项另作判断：只有上游及完整P输入相同，才在共同5工作米预算内比较省点率。上游变化后的省点率用于描述，不归因于DP本身变强。不同指标不临时加权成总分；缺少预先认可的可比标量目标时，也不计算一个容易误解的归一化regret。

\subsection{参考设置、数据分工与比较条件}
默认参数提供共同起点，而不是已知最佳答案。表\ref{tab:reference}列出参考设置。课件的95米来自骑行速度与采样间隔示例，starter另给400的初始化值；实验将两者作为比较对象，没有假设它们就是本数据的真实物理上限。

\begin{table}[htbp]\centering\small
\caption{基础处理的参考设置}\label{tab:reference}
\begin{tabularx}{\linewidth}{@{}YY@{}}\toprule
分段与过滤 & 方向与简化\\\midrule
时间差：30秒；相邻距离：400工作米 & 方向差：35度\\
最少点数：5；最短累计长度：65工作米 & DP容差：5工作米\\\bottomrule
\end{tabularx}
\end{table}

样本按用途划分，不能把重复使用过的数据重新称为独立检验。表\ref{tab:splits}给出各阶段的关系，阶段名称表示数据用途，而不是不同课程实验。

\begin{table}[htbp]\centering\small
\caption{数据的阶段用途与复用关系}\label{tab:splits}
\begin{tabularx}{\linewidth}{@{}L{31mm}rY@{}}\toprule
用途 & 记录数 & 参与什么判断\\\midrule
早期试运行 & 7 & 数学与执行回归、已暴露案例\\
示范记忆 & 60 & 建立并冻结经验，不从评估记录写回\\
初轮开发 & 120 & 参数、顺序与初步候选比较\\
阶段评估 & 120 & 检查已固定的单项、顺序和模式\\
组合开发 & 240 & 合并上述开发与已看过的阶段评估，不是另抽240条\\
新的选择数据 & 240 & 在固定短名单中取舍，不与最终确认混用\\
最终确认 & 600 & 最终设置固定后使用，不用于再调参\\\bottomrule
\end{tabularx}
\end{table}

四模式分别使用初轮开发与阶段评估内预先确定的24条共同记录。完整记录及已核实的关联组不跨分区，片段不能拆开制造独立样本。描述性分层依据原始跨度、时间和重复现象，不视为真实运动类别；分层样本的简单均值也不代表总体均值。

将选择和最终确认分开，是为了减轻反复尝试带来的选择偏差\cite{cawley}。组合开发中的旧评估数据已经暴露，后续不再赋予其未见测试的身份。最后的全量处理包含全部分区，属于生产统计，不是又一次独立确认。
